\chapter[Were Cave Paintings Just Decoration?]{Were Cave Paintings Merely Ancient Decoration?}
\section{A Red Stone}
Pick up something: a reddish-brown stone.

It fits comfortably in your palm. One face has worn flat like a blunt crayon. Draw it lightly across a grey wall and it leaves a clear red mark. This is ochre, an iron-bearing mineral found in many places worldwide.

Notice its oddity. Most prehistoric tools immediately reveal their uses: axes chop, bone needles sew, scrapers remove hide. This stone cannot feed, cut, or strike. Grinding takes time and effort; the red powder neither fills a stomach nor serves as medicine. Why would a hunter searching daily for food expend labour grinding a red stone to a sloping face?

It is no isolated example. Archaeologists have excavated thousands of ochre pieces worldwide: sharpened, powdered into shells, or engraved. South Africa's Blombos Cave yielded pieces over seventy thousand years old with orderly crosshatched incisions, over sixty millennia before farming or cities and roughly seventy before writing.

A useless stone, carefully ground, engraved, and carried. What for?

Follow it downward and you reach a great historical mystery: cave paintings. In French and Spanish caves, Sulawesi limestone caverns, Indian rock shelters, and Australian cliffs, people tens of thousands of years ago left bison, horses, handprints, and mysterious symbols. Their astonishing quality made the first scholars encountering them over a century ago refuse to believe primitive people could have made them.

Our question is simple: what were these paintings for? Decoration like posters and desk doodles, or a serious undertaking still undeciphered?

More important: with not one written word from their makers, what warrants our guesses, and when should guessing stop?

\section{An Afternoon in Altamira Cave}
Enter a scene: an afternoon in 1879, a cave in northern Spain's Cantabrian mountains.

The cave lies on Marcelino Sanz de Sautuola's land. This amateur archaeology enthusiast became fascinated after seeing Palaeolithic bones and figurines at an exhibition some years earlier. Today he brings young daughter Maria to excavate more ancient remains.

There is no electric light. Holding a lamp, he crouches near the entrance, scraping earth layer by layer with a small trowel, intent on bones and tools. Maria cannot stay still; with her own candle she explores the low spaces, looking around.

Then she looks up.

Above her, a herd of bison covers the low ceiling. Red, black, ochre; standing, lying, arching backs as if charging; the largest over a metre long. Artists used natural rock bulges to swell backs and bellies. Flickering candlelight sends shadows across bodies that almost breathe.

Later accounts say Maria called her father to look, something like: "Papa, look, bulls!"

He saw. He recognised great age: pigment penetrated stone, calcium deposits covered paintings, and entrance deposits contained Palaeolithic objects. In 1880 he self-published a booklet declaring Altamira's bison Palaeolithic work.

The most dramatic part followed. Almost all leading European prehistorians said impossible. Their reasoning sounded substantial: hunters struggling merely to survive could not create pictures admired even by modern painters. Some alleged forgery, others recent hunters' scribbles. Sautuola lacked academic credentials and evidence they would accept. At his 1888 death, mainstream scholarship still treated the bison as a joke.

Remember this afternoon and two pairs of eyes: a child's looking upward, academia's looking down. A child saw bison; experts refused to see their age. What decides who is right? That is our practice here.

\section{Lay the Question on the Table}
Clarify the conflict before choosing sides.

The decoration camp starts from plain intuition: humans naturally draw. We doodle during meetings, give textbook historical figures sunglasses and moustaches, choose patterned phone cases, paint walls. Drawing needs no reason, just as birds sing and cats groom. Rock walls are ancient notebook margins. Many paintings also occupy open, inhabited places, disorderly overlapping images resembling layered message walls rather than sacred plans. Decoration interferes least with ancient people: it invents no incomprehensible religion for them.

The opposing camp has strong cards too.

First, location. Altamira's bison cover a low ceiling; some French images require hundreds of metres of pitch-black passages narrow enough to demand crawling; others require scaffolding to reach. Who hangs decorations there? Home posters invite daily viewing. Risking deep darkness to paint must have a compelling reason.

Second, content. If decoration, why repeat so few subjects? Animals and symbols dominate; landscapes and plants are nearly absent; humans are rare and sketchy. A painter seeking mere prettiness should have a less peculiar menu.

Third, pigments and craft. Some minerals travelled dozens of kilometres. Some hand stencils required mouthfuls of pigment sprayed around a palm pressed to rock, making negative outlines and provoking coughing. Worthwhile for decoration?

Notice that the dispute concerns human motives, not artistic quality. One camp prefers the least mysterious explanation, the urge to draw. The other refuses to impose modern wall-decorating habits on a world without writing, where images might be supremely serious technology.

Choose a provisional side: first entering Altamira, would you think "How beautiful" or "What were they doing?"

\section{Who Speaks for the Cave's People?}
"Art for art's sake." After Altamira finally gained recognition early in the twentieth century, European critics offered this first explanation. Their era embraced art serving no purpose beyond beauty. Through that lens, cave painters became humanity's first artists, caves its first galleries, bison proof of an innate aesthetic impulse rooted tens of millennia earlier. Notice how pleasing this sounds: ancient hunters retrospectively become our earliest colleagues.

At almost the same time, others answered differently. Some hunter-gatherer peoples still painted rocks. Surely asking them beat guessing in a study?

Thus ethnography's voice entered. Anthropologists consulted Indigenous Australians and southern Africa's San. The answers discomforted art-for-art's-sake advocates: images had functions. Nineteenth-century records connect San paintings to rituals and rainmaking. Painting an eland meant capturing its power, not merely admiring it; some images showed people in ritual trance. Many Aboriginal Australian paintings tie ancestry, genealogy, and laws of land together, one image simultaneously map, legal provision, and family tree. Knowledge holders still explain particular patterns and their rules: some open to all, others restricted by status, others forbidden to reproduce.

Count the positions and their social locations.

First, nineteenth-century academics judged prehistoric people too primitive to paint: prejudice, not explanation, later falsified. Second, amateur Sautuola brought eyes and evidence without credentials and went unheard in life. Third, aesthetic critics imposed their era's slogan on distant prehistory. Fourth, ethnographers and Indigenous knowledge holders illuminated lost images through living traditions. Their own caution matters equally: San motives cannot automatically become French cave painters' motives across half a world and over ten thousand years.

Another detail shows how choosing whom to hear changes judgement. A famous Helan Mountain rock image, the "Sun God", presents a face surrounded by radiating lines. Tourist guides confidently identify sun worship: surely rays mean sun. Researchers familiar with northern pastoral traditions note similar radiating ritual headdresses. Calling it the sun may impose our conclusion. Both readings remain alive; the wall stays silent. The small example contains the entire problem: everybody sees the image, but its meaning depends on whose knowledge they bring.

The most important lesson is that those nearest the primary material, Indigenous knowledge holders, long counted as research subjects rather than interpreters. Once scholars listened seriously, the field's understanding reversed. Whom we listen to is itself part of method.

\section[Thinking Bridge: Remains to Behaviour]{A Thinking Bridge: Infer Behaviour from What Remains}
Learn our portable tool: how archaeologists reconstruct vanished behaviour from silent remnants.

Start with an emptied, unswept classroom after school. Chalk dust suggests lessons; sweet wrappers inside back-row desks suggest eating in class; broken chalk suggests someone snapped it, perhaps teacher or playful pupil. Confidence differs: lessons are nearly certain; eating likely; who and why less secure. Whether the eater was happy is beyond the remnants' speech.

Archaeologists facing caves are in precisely this position. They divide inference into three layers with different methods and certainty.

First, facts: what things are and their age. Hard tools help. Radiocarbon dates charcoal pigment; uranium--thorium dating measures thin calcium-carbonate films above images. Calcite over a Sulawesi wild-pig painting dates to at least forty-five thousand years ago, so the painting must be older. This layer is firmest.

Second, behaviour: what happened. Analogy and experiment help. To investigate hand stencils, experimental archaeologists actually mouth ochre slurry and spray it around hands on stone. Possible, but choking, requiring repeated breaths. To estimate painting time, people copy images by torchlight in simulated caves, recording time and pigment. To test stone lamps, they reproduce excavated forms, burn animal fat, and assess flame stability and smoke. Experiments cannot prove ancient people used that method, but establish possibility and separate it from impossibility.

Third, meaning: why. Tools grow blunter. Ethnographic analogy suggests possibilities, including ritual uses of images, but analogy inspires rather than testifies. Statistical patterns in location, orientation, and combinations reveal regularities, not necessarily meanings.

The principle: certainty declines layer by layer, and honest people label each layer's quality. Hearing prehistoric claims such as shamanic authorship or rainmaking, first ask which layer supports them and what lies beneath.

Here is our boundary: never pretend to know intentions the makers left unrecorded. Tools bring us closer to truth, but every mark on their scale should state our confidence.

\section{An Essay That Made Scholars Admit Error}
Read primary material, not prehistoric writing, which does not exist here, but something equally important: a scholar's public admission of error.

In 1902, leading French prehistorian Émile Cartailhac, once prominent in rejecting Altamira, examined newly discovered paintings in L'Anthropologie. His subtitle openly confessed:

\begin{quote}
"Mea culpa d'un sceptique"---a sceptic's "my mistake"
\end{quote}
The argument: I was wrong. New evidence, particularly paintings covered by stalagmitic layers and associated with Palaeolithic strata, removes reasonable doubt. Altamira and these other paintings are genuinely Palaeolithic.

The essay deserves close reading not merely for proving ancient drawing, but for showing a discipline changing its mind. Rhetoric did not turn scholarship; the evidence changed shape. One Altamira could be isolated, forged, coincidental. Successive southern French discoveries, with millennia-old calcite over paintings and datable Palaeolithic remains nearby, imposed increasing explanatory burdens on forgery until it became more absurd than antiquity. The burden shifted; the discipline turned.

Do not miss another detail: Sautuola had been dead fourteen years. The apology arrived too late for him. Evidence may win without winning fast enough to do justice to individuals.

Dating now extends the story further. Here are widely cited approximate dates with reliable orders of magnitude:

\begin{tabularx}{\linewidth}{llX}
\toprule
Site & Remains & Approximate age \\
\midrule
Blombos Cave, South Africa & Crosshatched ochre, shell beads & Over 70,000 years \\
Sulawesi, Indonesia & Wild-pig images, handprints & At least 45,000 years \\
Chauvet Cave, France & Lions, rhinoceroses, horses & Over 30,000 years \\
Lascaux Cave, France & Hundreds of horses, bison, deer & About 17,000 years \\
Altamira, Spain & Ceiling bison & Over 10,000 to over 30,000 years; painted repeatedly \\
\bottomrule
\end{tabularx}
Our strongest evidence hand. What does it prove? Humans have made rock images across continents for at least forty-five thousand years, while the earliest generally accepted Sumerian and Egyptian writing is only over five thousand years old. Image history is several times longer. The table does not establish why. Every date belongs to facts; meanings remain contested interpretations.

Another cave makes location's extremes vivid. Cosquer Cave on France's Mediterranean coast now opens dozens of metres below sea level. Divers traverse a dark underwater tunnel to chambers containing horses, bison, and birds no longer found in those waters. During painting, the ice age continued, sea level was far lower, and the entrance stood on land. The images themselves witness climate change. Consider the weight of this: environmental transformation submerged the entrance, while paintings waited in place for tens of millennia until a modern diver relit them.

Now enter the interpretive battle.

\section{Five Explanations on the Table}
First distinguish four things; everything following depends on it.

What happened: people across continents left animals, hands, and symbols on rock tens of thousands of years ago, established evidence. Why: interpretation, where five accounts will compete. Participants' understanding: no written record, leaving this layer nearly empty. Honestly acknowledge that other hunter-gatherers' accounts offer analogy, not testimony. Our evaluation: calling these great artistic treasures describes us, not their makers.

Now begin.

\textbf{Explanation 1: Decoration and play.} Give it a full hearing; mocking it for reducing great art to doodles is unfair. Three reasons support it. First, drawing is universally observable, from toddlers to adults and desks to subway advertisements. Ordinary motives make economical scientific hypotheses. Second, many outdoor, residential images overlap freely without obvious order. Guangxi's red Zuojiang Huashan images occupy riverside cliffs; Yinshan and Helan images spread across desert rocks. They resemble layered public messages more than hidden icons. Third, sacredness and ritual are archaeologists' easily abused bins for anything unexplained. Decoration protects a discipline: do not conceal ignorance behind mystery. Its weakness is familiar: casual doodling poorly explains dangerous depths and heights, or travelling dozens of kilometres for a mouthful of pigment.

\textbf{Explanation 2: Hunting magic.} Popular early in the twentieth century: depicting prey enables its subjugation; a painted spear strike ensures tomorrow's real one. Apparent evidence includes spear- or trap-like marks on animals and finger-dotted surroundings suggesting spells. Elegant, specific, testable---and tested. Archaeologists counted food remains at sites including Lascaux: reindeer dominated meals, horses and bison the walls. Painted animals were not those eaten. Why not paint more reindeer to secure prey? This easily missed counterexample shows seemingly supportive evidence unravel when kitchen waste is checked. Advocates can rescue it with important rather than commonly eaten animals, but the theory begins accumulating patches.

\textbf{Explanation 3: Ritual and shamanism.} Among today's most influential. Fever, oxygen deprivation, meditation, and trance can generate geometric visual phenomena: zigzags, spirals, grids, luminous points. These entoptic phenomena arise from nervous systems shared across ages. Long-neglected cave symbols frequently have such forms. Dark, sensory-isolated depths could therefore encourage trance; paintings record ritual visions; palms on rock touch a world beyond it. This gives abstract symbols a role and matches San ethnographic records. Its weakness: universal capacity for visions does not mean universal painting of them. One neurological fact stretches too far if applied to every image on every continent across millennia.

\textbf{Explanation 4: Structure and symbolism.} French scholar André Leroi-Gourhan statistically examined animal placement and groupings, proposing ordered distributions: horses and bison in particular positions might represent genders or groups, making the cave a complete symbolic arrangement. Its advantage is reading the cave as a book, not leafing through a picture album. Its weakness: small samples; later recounts found less stable patterns.

\textbf{Explanation 5: Memory and identity.} Plain and easily underestimated: images externalise memory. Without writing they store genealogies, mark territory, record routes, and declare who we are and where we belong. Aboriginal Australian traditions show these functions: one pattern can combine deed and genealogy. Handprints especially resemble signatures: I came, I existed. Like decoration, this avoids excessive mystery, yet explains depth as a solemn anchor for identity and memory. Its weakness: near-unfalsifiability. Almost any image could be memory.

Five explanations grasp differently, without complete mutual exclusion. One wall may begin as doodling and become sacred; one handprint may sign and contact spirits. The methodological lesson is that each grasps part of the evidence and declares it the whole. Whenever a theory explains everything, remember Lascaux's reindeer bones: evidence's drawer is always larger than theory.

\section[Further Challenge: Beauty Before Writing]{A Deeper Challenge: Why Is Beauty Older Than Writing?}
Go deeper through aesthetics, the study of beauty itself. Why might aesthetic experience predate cities, writing, even farming by so long?

First perform conceptual surgery. In the eighteenth-century Critique of Judgement, Kant famously describes aesthetic pleasure as disinterested. Salivating over an apple is appetite; wanting a house is desire for possession. Sunset pleasure without wanting to eat or own it is aesthetic. Beauty makes us pause before something useless.

Return to the red stone. Grinding ochre takes hours; pigment cannot feed us; a painting hundreds of metres into darkness cannot light our homeward path. These acts fall outside usefulness. Kant's knife makes them candidates for disinterested activity. A Sulawesi painter forty-five thousand years ago spends potential foraging time mixing pigment deep underground for a wild pig. If that is not pausing for beauty, what is?

Do not rush. One honest obstacle remains. Costly investment in apparently useless things is observable; aesthetic pleasure is an invisible inner state. Reverent ritual can be equally solemn and disinterested. Carefully stated, paintings establish ancient willingness and ability to pay high costs for something beyond staying alive. They strongly suggest aesthetic experience but cannot prove the painter felt precisely our sunset pleasure. Keeping that distinction avoids inventing unrecorded intention.

Consider older, easily misjudged evidence. A natural pebble from South Africa's Makapansgat was carried several kilometres to a camp around three million years ago by australopithecines, hominins predating erectus. Unworked, it naturally resembles a face: two eye-like hollows and a mouthlike crack. Why carry an inedible stone home? The most reasonable guess is that its facial resemblance was noticed and found interesting.

One stone is no artwork; it was not even worked by a finger. Yet its suggestion runs deep: being moved by shapes may predate our human species. Symmetry recognition, attraction to bright colours, and seeing resemblance lie deep in nervous systems, long before anyone coined art.

The mystery's final layer therefore changes the question. Perhaps asking whether ancient people had art is mistaken. Art as concept, institution, and market is indeed recent. Aesthetic experience is ancient: stopping before useless things, being caught, wanting to expend effort. Writing and cities are over five thousand years old, farming over ten thousand. The impulse to press red hands onto rock is at least nine times writing's age and four times farming's.

Beauty is not civilisation's fruit. Perhaps the reverse is nearer truth: civilisation descends from beauty.

\section{You Too Leave Handprints on Walls}
This forty-five-thousand-year-old question sits beside you.

You continually leave handprints: textbook doodles, desk scratches, phone-case patterns, social-media covers, game skins, screen comments. The urge to declare presence never left us. Schoolbooks recount young Lu Xun carving zao, "early", into his desk as a reminder. Over a century later, that desk sits in a museum with countless visitors inspecting the mark. Time can retrospectively make a child's graffiti a relic. Yesterday's vandalism becomes tomorrow's heritage through time and narration.

Today's world stages every explanation. Street graffiti reenacts the cave dispute: authorities see dirt, artists expression, residents perhaps simply a new cat on the wall downstairs. Famous graffiti fetches astronomical auction prices; anonymous teenagers' work is painted over. Art versus vandalism depends not merely on images but names and locations. Is this not identity theory live?

On-screen comments and memes place words and images directly over pictures and others' speech, almost a new act in media history performing the oldest function: proof of presence, "I am here too". Overlapping deep-cave hands and crowded screen comments perform one act across tens of millennia.

Ritual and identity explanations revive online too. Personal taglines, avatar decorations, and profile backgrounds are digital totems and territorial marks. Your carefully tended profile is your cave's depths, arranged as an exhibition of identity, whose visitors matter.

Pause longest at the new variable: algorithms, not pigments and torches, now decide which images billions see. They allocate wall space by viewing time, an ancient system driven entirely by the desire to be seen. Once visibility required a torchbearer entering darkness; now a machine calculates which picture stops a scrolling finger. Only the torch changes: people still willingly pause before images.

\section[Your Turn: What to Leave in the Dark]{Your Turn: What Is Worth Leaving in the Darkness?}
Return to that Altamira afternoon.

When Maria looked up, over ten thousand years separated her from the painter, with no words crossing the gap. She knew neither name, motive, meal after painting, nor any upward glance and remark to companions. All that remains submerged forever. She had only the red-and-black, almost breathing herd.

That leaves our final question to you:

If intention can never be established, if humanity knew it could never answer why they painted, could we still discuss the pictures' meaning? Or should complete honesty say only "We do not know" and fall silent?

Count your cards again. For discussion: interpretation is not invention; all five accounts face evidence, and reindeer bones can eliminate bad theories. Missing details is not total ignorance: we know costs, locations, shared perception. For silence: confident explanations once became consensus, then another cave overturned them. Our prehistoric stories often narrate ourselves: aesthetic autonomy describes twentieth-century Europe; hunting magic echoes nineteenth-century ideas of primitive religion.

A personal version: tomorrow you must leave a mark somewhere unvisited for a long time, addressed to its eventual discoverer. Animal, handprint, or words? How much would readers tens of thousands of years hence guess correctly?

No standard answer exists. One thing is certain: by taking the question to heart, you join the chain. Between hands pressed in ancient darkness and your eyes on this page lies not an abyss but an unbroken wish to leave something behind.
